\documentclass{amsart}
% For dealing with references we use the comment environment
\usepackage{verbatim}
\newenvironment{reference}{\comment}{\endcomment}
%\newenvironment{reference}{}{}
\newenvironment{slogan}{\comment}{\endcomment}
\newenvironment{history}{\comment}{\endcomment}

% For commutative diagrams we use Xy-pic
\usepackage[all]{xy}

% We use 2cell for 2-commutative diagrams.
\xyoption{2cell}
\UseAllTwocells

% We use multicol for the list of chapters between chapters
\usepackage{multicol}

% This is generally recommended for better output
\usepackage{lmodern}
\usepackage[T1]{fontenc}

% For cross-file-references

% Package for hypertext links:
\usepackage{hyperref}

% For any local file, say "hello.tex" you want to link to please

% Theorem environments.
%
\theoremstyle{plain}
\newtheorem{theorem}[subsection]{Theorem}
\newtheorem{proposition}[subsection]{Proposition}
\newtheorem{lemma}[subsection]{Lemma}

\theoremstyle{definition}
\newtheorem{definition}[subsection]{Definition}
\newtheorem{example}[subsection]{Example}
\newtheorem{exercise}[subsection]{Exercise}
\newtheorem{situation}[subsection]{Situation}

\theoremstyle{remark}
\newtheorem{remark}[subsection]{Remark}
\newtheorem{remarks}[subsection]{Remarks}

\numberwithin{equation}{subsection}

% Macros
%
\def\lim{\mathop{\mathrm{lim}}\nolimits}
\def\colim{\mathop{\mathrm{colim}}\nolimits}
\def\Spec{\mathop{\mathrm{Spec}}}
\def\Hom{\mathop{\mathrm{Hom}}\nolimits}
\def\Ext{\mathop{\mathrm{Ext}}\nolimits}
\def\SheafHom{\mathop{\mathcal{H}\!\mathit{om}}\nolimits}
\def\SheafExt{\mathop{\mathcal{E}\!\mathit{xt}}\nolimits}
\def\Sch{\mathit{Sch}}
\def\Mor{\mathop{\mathrm{Mor}}\nolimits}
\def\Ob{\mathop{\mathrm{Ob}}\nolimits}
\def\Sh{\mathop{\mathit{Sh}}\nolimits}
\def\NL{\mathop{N\!L}\nolimits}
\def\CH{\mathop{\mathrm{CH}}\nolimits}
\def\proetale{{pro\text{-}\acute{e}tale}}
\def\etale{{\acute{e}tale}}
\def\QCoh{\mathit{QCoh}}
\def\Ker{\mathop{\mathrm{Ker}}}
\def\Im{\mathop{\mathrm{Im}}}
\def\Coker{\mathop{\mathrm{Coker}}}
\def\Coim{\mathop{\mathrm{Coim}}}

% Boxtimes
%
\DeclareMathSymbol{\boxtimes}{\mathbin}{AMSa}{"02}

%
% Macros for moduli stacks/spaces
%
\def\QCohstack{\mathcal{QC}\!\mathit{oh}}
\def\Cohstack{\mathcal{C}\!\mathit{oh}}
\def\Spacesstack{\mathcal{S}\!\mathit{paces}}
\def\Quotfunctor{\mathrm{Quot}}
\def\Hilbfunctor{\mathrm{Hilb}}
\def\Curvesstack{\mathcal{C}\!\mathit{urves}}
\def\Polarizedstack{\mathcal{P}\!\mathit{olarized}}
\def\Complexesstack{\mathcal{C}\!\mathit{omplexes}}
% \Pic is the operator that assigns to X its picard group, usage \Pic(X)
% \Picardstack_{X/B} denotes the Picard stack of X over B
% \Picardfunctor_{X/B} denotes the Picard functor of X over B
\def\Pic{\mathop{\mathrm{Pic}}\nolimits}
\def\Picardstack{\mathcal{P}\!\mathit{ic}}
\def\Picardfunctor{\mathrm{Pic}}
\def\Deformationcategory{\mathcal{D}\!\mathit{ef}}

\setlength{\emergencystretch}{3em}
\begin{document}
\title[Stacks Project, Tag 05RG]{316 --- Verdier multiplicative systems associated to triangulated subcategories}
\maketitle
\noindent Copyright (C) 2005--2025 Johan de Jong. Stacks Project authors. Source: \href{https://stacks.math.columbia.edu/tag/05RG}{Tag 05RG}.

\noindent Editorial difficulty note (not author text). Difficulty H2: The proof verifies the Ore and triangle-compatibility conditions using octahedra and the nine-triangle construction. The converse saturation implication requires an additional cone diagram, a zero composite and splitting to recover the individual cones. This establishes the localization machinery and the precise relation between saturation of arrows and subcategories..

\noindent Editorial provenance note (not author text). History: excerpt from revision \texttt{a0444\allowbreak 6e57e\allowbreak c1fbc\allowbreak 252a8\allowbreak 71afc\allowbreak ec775\allowbreak 2fb28\allowbreak 07b14}, derived.tex, lines 1908--2063. Reference presentation was adapted; original-author context and core support blocks were retained or added. The mathematical wording is unchanged. Licensed under GNU FDL 1.2 or later; no invariant sections or cover texts. See ../licenses/COPYING. Context blocks reproduce author definitions and supporting results; they are not additional counted problems.

\begin{definition}
\label{definition-localization}
Let $\mathcal{D}$ be a pre-triangulated category. We say a multiplicative
system $S$ is {\it compatible with the triangulated structure} if
the following two conditions hold:
\begin{enumerate}
\item[MS5] For a morphism $f$ of $\mathcal{D}$ we have
$f \in S \Leftrightarrow f[1] \in S$\footnote{See Remark \href{https://stacks.math.columbia.edu/tag/0H30}{remark MS5 (Tag 0H30)}.}.
\item[MS6] Given a solid commutative square
$$
\xymatrix{
X \ar[r] \ar[d]^s &
Y \ar[r] \ar[d]^{s'} &
Z \ar[r] \ar@{-->}[d] &
X[1] \ar[d]^{s[1]} \\
X' \ar[r] &
Y' \ar[r] &
Z' \ar[r] &
X'[1]
}
$$
whose rows are distinguished triangles with $s, s' \in S$
there exists a morphism $s'' : Z \to Z'$ in $S$ such that
$(s, s', s'')$ is a morphism of triangles.
\end{enumerate}
\end{definition}

\begin{definition}
\label{definition-triangulated-subcategory}
Let $(\mathcal{D}, [\ ], \mathcal{T})$ be a pre-triangulated category.
A {\it pre-triangulated subcategory}\footnote{This definition may be
nonstandard. If $\mathcal{D}'$ is a full subcategory then $\mathcal{T}'$
is the intersection of the set of triangles in $\mathcal{D}'$ with
$\mathcal{T}$, see
Lemma \href{https://stacks.math.columbia.edu/tag/05QX}{lemma triangulated subcategory (Tag 05QX)}.
In this case we drop $\mathcal{T}'$ from the notation.}
is a pair $(\mathcal{D}', \mathcal{T}')$ such that
\begin{enumerate}
\item $\mathcal{D}'$ is an additive subcategory of $\mathcal{D}$
which is preserved under $[1]$ and such that
$[1] : \mathcal{D}' \to \mathcal{D}'$ is an auto-equivalence,
\item $\mathcal{T}' \subset \mathcal{T}$ is a subset such that for every
$(X, Y, Z, f, g, h) \in \mathcal{T}'$ we have
$X, Y, Z \in \Ob(\mathcal{D}')$ and
$f, g, h \in \text{Arrows}(\mathcal{D}')$, and
\item $(\mathcal{D}', [\ ], \mathcal{T}')$ is a pre-triangulated
category.
\end{enumerate}
If $\mathcal{D}$ is a triangulated category, then we say
$(\mathcal{D}', \mathcal{T}')$ is a {\it triangulated subcategory} if
it is a pre-triangulated subcategory and
$(\mathcal{D}', [\ ], \mathcal{T}')$ is a triangulated category.
\end{definition}

\begin{definition}
\label{definition-saturated}
Let $\mathcal{D}$ be a pre-triangulated category. We say a full
pre-triangulated subcategory $\mathcal{D}'$ of $\mathcal{D}$ is
{\it saturated} if whenever $X \oplus Y$ is isomorphic to an object
of $\mathcal{D}'$ then both $X$ and $Y$ are isomorphic to objects
of $\mathcal{D}'$.
\end{definition}

\begin{lemma}
\label{lemma-construct-multiplicative-system}
Let $\mathcal{D}$ be a triangulated category.
Let $\mathcal{D}' \subset \mathcal{D}$ be a full triangulated
subcategory. Set
\begin{equation}
\label{equation-multiplicative-system}
S =
\left\{
\begin{matrix}
f \in \text{Arrows}(\mathcal{D})
\text{ such that there exists a distinguished triangle }\\
(X, Y, Z, f, g, h) \text{ of }\mathcal{D}\text{ with }
Z\text{ isomorphic to an object of }\mathcal{D}'
\end{matrix}
\right\}
\end{equation}
Then $S$ is a multiplicative system compatible with the triangulated
structure on $\mathcal{D}$. In this situation the following are equivalent
\begin{enumerate}
\item $S$ is a saturated multiplicative system,
\item $\mathcal{D}'$ is a saturated triangulated subcategory.
\end{enumerate}
\end{lemma}

\begin{proof}
To prove the first assertion we have to prove that
MS1, MS2, MS3 and MS5, MS6 hold.

\medskip\noindent
Proof of MS1. It is clear that identities are in $S$ because
$(X, X, 0, 1, 0, 0)$ is distinguished for every object $X$ of $\mathcal{D}$
and because $0$ is an object of $\mathcal{D}'$. Let $f : X \to Y$
and $g : Y \to Z$ be composable morphisms contained in $S$.
Choose distinguished triangles $(X, Y, Q_1, f, p_1, d_1)$,
$(X, Z, Q_2, g \circ f, p_2, d_2)$, and $(Y, Z, Q_3, g, p_3, d_3)$.
By assumption we know that $Q_1$ and $Q_3$ are isomorphic to objects
of $\mathcal{D}'$. By TR4 we know there exists a distinguished
triangle $(Q_1, Q_2, Q_3, a, b, c)$. Since $\mathcal{D}'$ is a
triangulated subcategory we conclude that $Q_2$ is isomorphic to
an object of $\mathcal{D}'$. Hence $g \circ f \in S$.

\medskip\noindent
Proof of MS3. Let $a : X \to Y$ be a morphism and let $t : Z \to X$ be
an element of $S$ such that $a \circ t = 0$. To prove LMS3 it suffices to
find an $s \in S$ such that $s \circ a = 0$, compare with the proof of
Lemma \href{https://stacks.math.columbia.edu/tag/05R4}{lemma triangle functor localize (Tag 05R4)}. Choose a distinguished
triangle $(Z, X, Q, t, g, h)$ using TR1 and TR2. Since $a \circ t = 0$
we see by
Lemma \href{https://stacks.math.columbia.edu/tag/0149}{lemma representable homological (Tag 0149)}
there exists a morphism $i : Q \to Y$ such that $i \circ g = a$.
Finally, using TR1 again we can choose a triangle
$(Q, Y, W, i, s, k)$. Here is a picture
$$
\xymatrix{
Z \ar[r]_t & X \ar[r]_g \ar[d]^1 & Q \ar[r] \ar[d]^i & Z[1] \\
& X \ar[r]_a & Y \ar[d]^s \\
& & W
}
$$
Since $t \in S$ we see that $Q$ is isomorphic to an object of $\mathcal{D}'$.
Hence $s \in S$. Finally, $s \circ a = s \circ i \circ g = 0$ as
$s \circ i = 0$ by Lemma \href{https://stacks.math.columbia.edu/tag/0146}{lemma composition zero (Tag 0146)}.
We conclude that LMS3 holds.
The proof of RMS3 is dual.

\medskip\noindent
Proof of MS5. Follows as distinguished triangles and $\mathcal{D}'$
are stable under translations

\medskip\noindent
Proof of MS6. Suppose given a commutative diagram
$$
\xymatrix{
X \ar[r] \ar[d]^s &
Y \ar[d]^{s'} \\
X' \ar[r] &
Y'
}
$$
with $s, s' \in S$. By
Proposition \href{https://stacks.math.columbia.edu/tag/05R0}{proposition 9 (Tag 05R0)}
we can extend this to a nine square diagram. As $s, s'$ are elements of $S$
we see that $X'', Y''$ are isomorphic to objects of $\mathcal{D}'$.
Since $\mathcal{D}'$ is a full triangulated subcategory we see that
$Z''$ is also isomorphic to an object of $\mathcal{D}'$.
Whence the morphism $Z \to Z'$
is an element of $S$. This proves MS6.

\medskip\noindent
MS2 is a formal consequence of MS1, MS5, and MS6, see
Lemma \href{https://stacks.math.columbia.edu/tag/05R3}{lemma localization conditions (Tag 05R3)}.
This finishes the proof of the first assertion of the lemma.

\medskip\noindent
Let's assume that $S$ is saturated. (In the following we will use
rotation of distinguished triangles without further mention.)
Let $X \oplus Y$ be an object isomorphic to an object of $\mathcal{D}'$.
Consider the morphism $f : 0 \to X$. The composition
$0 \to X \to X \oplus Y$ is an element
of $S$ as $(0, X \oplus Y, X \oplus Y, 0, 1, 0)$ is a distinguished
triangle. The composition $Y[-1] \to 0 \to X$ is an element of $S$
as $(X, X \oplus Y, Y, (1, 0), (0, 1), 0)$ is a distinguished triangle, see
Lemma \href{https://stacks.math.columbia.edu/tag/05QT}{lemma split (Tag 05QT)}.
Hence $0 \to X$ is an element of $S$ (as $S$ is saturated).
Thus $X$ is isomorphic to an object of $\mathcal{D}'$ as desired.

\medskip\noindent
Finally, assume $\mathcal{D}'$ is a saturated triangulated subcategory.
Let
$$
W \xrightarrow{h}
X \xrightarrow{g}
Y \xrightarrow{f} Z
$$
be composable morphisms of $\mathcal{D}$ such that $fg, gh \in S$.
We will build up a picture of objects as in the diagram below.
$$
\xymatrix{
 & &
Q_{12} \ar[rd] & &
Q_{23} \ar[rd] \\
 &
Q_1 \ar[ld]_{\! + \! 1} \ar[ru] & &
Q_2 \ar[ld]_{\! + \! 1} \ar[ll]_{\! + \! 1} \ar[ru] & &
Q_3 \ar[ld]_{\! + \! 1} \ar[ll]_{\! + \! 1} \\
W \ar[rr] & &
X \ar[lu] \ar[rr] & &
Y \ar[lu] \ar[rr] & &
Z \ar[lu]
}
$$
First choose distinguished triangles
$(W, X, Q_1)$, $(X, Y, Q_2)$, $(Y, Z, Q_3)$ $(W, Y, Q_{12})$, and
$(X, Z, Q_{23})$. Denote $s : Q_2 \to Q_1[1]$ the composition
$Q_2 \to X[1] \to Q_1[1]$. Denote $t : Q_3 \to Q_2[1]$ the
composition $Q_3 \to Y[1] \to Q_2[1]$.
By TR4 applied to the composition $W \to X \to Y$
and the composition $X \to Y \to Z$ there exist
distinguished triangles $(Q_1, Q_{12}, Q_2)$ and $(Q_2, Q_{23}, Q_3)$
which use the morphisms $s$ and $t$.
The objects $Q_{12}$ and $Q_{23}$ are isomorphic to objects of
$\mathcal{D}'$ as $W \to Y$ and $X \to Z$ are assumed in $S$.
Hence also $s[1]t$ is an element of $S$ as $S$ is closed under compositions
and shifts.
Note that $s[1]t = 0$ as $Y[1] \to Q_2[1] \to X[2]$ is zero, see
Lemma \href{https://stacks.math.columbia.edu/tag/0146}{lemma composition zero (Tag 0146)}.
Hence $Q_3[1] \oplus Q_1[2]$ is isomorphic to an object of
$\mathcal{D}'$, see Lemma \href{https://stacks.math.columbia.edu/tag/05QT}{lemma split (Tag 05QT)}.
By assumption on $\mathcal{D}'$ we conclude that $Q_3$ and $Q_1$ are isomorphic
to objects of $\mathcal{D}'$. Looking at the distinguished triangle
$(Q_1, Q_{12}, Q_2)$ we conclude that $Q_2$ is also isomorphic to
an object of $\mathcal{D}'$. Looking at the distinguished triangle
$(X, Y, Q_2)$ we finally conclude that $g \in S$. (It is also
follows that $h, f \in S$, but we don't need this.)
\end{proof}
\end{document}
